\chapter{Model Risk and Validation}\label{ch:rc:model-risk-and-validation}

On 30 January 2012 JPMorgan's model review group authorised a new value-at-risk model for the
synthetic credit portfolio of the bank's Chief Investment Office. The old model had been judged
too conservative; the new one cut the office's VaR by half, from 132 to 66 million dollars, and the
firm-wide VaR limit, breached for days, was no longer exceeded. Months later, after losses on the portfolio had grown to billions of dollars, the
bank's own task force found an operational error in the model's spreadsheets: after subtracting
an old hazard rate or correlation from a new one, the spreadsheet divided by their sum instead of
their average, which likely muted volatility by a factor of two and lowered the VaR. The model had
been approved; nobody had checked that it computed what its documentation said. This chapter is
about the discipline that does: identifying models, tiering them by the damage they can do,
validating them against independent benchmarks and against outcomes, and governing the results.

\section{What model risk is and how it enters}

\begin{definition}[Model risk]\label{def:rc:model-risk-and-validation:risk}
\emph{Model risk}\index{model risk} is the potential for adverse consequences from decisions
based on a model's output: financial loss, misstatement, or poor risk decisions. It enters through
wrong theory or assumptions, errors of implementation and data, and use of a sound model outside
the conditions it was built for.
\end{definition}

\begin{definition}[Model risk management]\label{def:rc:model-risk-and-validation:mrm}
\emph{Model risk management}\index{model risk management} is the framework that identifies a firm's
models, assesses and limits the risk they carry, and assigns responsibility for development, use,
validation and oversight across their life.
\end{definition}

The three sources call for different checks. Wrong theory is caught by conceptual review and by
comparing against other models; implementation errors by independent re-implementation and by
tests on cases with known answers; misuse by monitoring the model's outcomes and its inputs.

\section{Supervisory expectations}

\begin{definition}[Effective challenge]\label{def:rc:model-risk-and-validation:challenge}
\emph{Effective challenge}\index{effective challenge} is the critical analysis of a model by objective
people with the expertise to find its weaknesses, the independence to report them, and the standing
to have them fixed.
\end{definition}

\begin{definition}[Three lines of defence]\label{def:rc:model-risk-and-validation:lines}
The \emph{three lines of defence}\index{three lines of defence} assign \omterm{def:rc:model-risk-and-validation:risk}{model risk} to the model's
owners and users (first line: development and use), to an independent risk function (second line:
validation and oversight) and to internal audit (third line: assurance that the framework works).
\end{definition}

\begin{dated}{2026-09}{Model risk guidance}\label{dat:rc:model-risk-and-validation:rules}
In the United States, the Federal Reserve, the OCC and the FDIC issued revised guidance on \omterm{def:rc:model-risk-and-validation:mrm}{model risk
management} on 17 April 2026 (SR 26-2), superseding the guidance of April 2011 (SR 11-7). It is most
relevant to banking organisations with over 30 billion dollars of assets; it defines a model as a
complex quantitative method applying statistical, economic or financial theories to turn input data
into estimates, excluding simple arithmetic such as spreadsheet calculations and deterministic rule-based
processes, and it leaves generative and agentic AI outside its scope. In the United Kingdom, the PRA's
supervisory statement SS1/23 (published May 2023) has applied since 17 May 2024 to banks with internal
models for regulatory capital, through five principles: identification and classification, governance,
development and use, independent validation, and risk mitigants. In the euro area, the ECB revised its
guide to internal models on 28 July 2025, adding expectations for machine-learning techniques.
\end{dated}

\section{The validation report}

\begin{definition}[Model validation]\label{def:rc:model-risk-and-validation:validation}
\emph{Model validation}\index{model validation} is the set of activities that verify a model performs
as intended: conceptual soundness, verification of the implementation, benchmarking, and \omterm{def:rc:model-risk-and-validation:bench}{outcomes
analysis}, with findings, their severity, and conditions of use.
\end{definition}

\begin{definition}[Benchmark model, outcomes analysis]\label{def:rc:model-risk-and-validation:bench}
A \emph{benchmark model}\index{benchmark model} is an independent implementation, of the same
model or of an alternative, against which a model's outputs are compared. \emph{Outcomes
analysis}\index{outcomes analysis} compares a model's outputs with the realised outcomes they
predicted: backtesting of VaR and margin, P\&L attribution, default rates against PDs.
\end{definition}

\begin{method}[Benchmarking over a grid]\label{met:rc:model-risk-and-validation:grid}
Choose the parameter ranges the model will meet in use. Evaluate the model and the benchmark on
every combination. Set tolerances in business terms (an absolute error in basis points of notional,
and a relative error). Record every point, summarise failures by parameter, and investigate the
pattern, not only the worst case.
\end{method}

\omcode{code/firm/modelval/firm_modelval.py}{57}{74}{A benchmarking harness: every grid point recorded, with absolute and relative tolerances, and the failures summarised.}\label{lst:rc:model-risk-and-validation:bench}

\begin{example}[Validating chapter~7's tree]\label{ex:rc:model-risk-and-validation:tree}
Chapter~7's Hull--White trinomial tree prices European swaptions; the Jamshidian closed form of the
same model is the benchmark. Over 216 points (mean reversion 1\%, 5\%, 20\%; volatility 60 and 120
basis points; expiries 1, 5, 10 years; tenors 5 and 10; strikes at the money and $\pm1\%$; quarterly and
monthly steps), with a tolerance of half a basis point of notional or 1\%: with quarterly steps 54 of 108
points fail and the worst error is 15.8 basis points of notional; with monthly steps 23 fail and the worst
is 3.2. The failures concentrate at short expiries, 14 of the 23 at one year
(\cref{fig:rc:model-risk-and-validation:validation}): a one-year-into-five-year at-the-money payer
(mean reversion 1\%, volatility 60 basis points) is worth 104.81 basis points in closed form, 103.07 in the
monthly tree and 98.55 in the quarterly one. The finding: the tree needs finer steps before a near
exercise date, and a condition of use until it has them.
\end{example}

\begin{omfigure}
\begin{tikzpicture}
\begin{axis}[width=11.5cm,height=5.2cm,ybar,bar width=14pt,
  symbolic x coords={1y,5y,10y},xtick=data,enlarge x limits=0.3,
  xlabel={swaption expiry},ylabel={largest error (bp of notional)},
  tick label style={font=\small},label style={font=\small},
  ymin=0,ymax=18,grid=major,grid style={gray!25},
  nodes near coords,every node near coord/.append style={font=\footnotesize,/pgf/number format/fixed,/pgf/number format/precision=1},
  legend columns=2,legend style={font=\footnotesize,at={(0.5,-0.32)},anchor=north,draw=none,fill=none,/tikz/every even column/.append style={column sep=8pt}},
  table/col sep=comma]
\addplot[fill=omThm!60,draw=omThm] table[x=expiry,y=err_q]{figdata/rates-credit-risk/26-model-risk-and-validation/validation.csv};
\addplot[fill=omDef!60,draw=omDef] table[x=expiry,y=err_m]{figdata/rates-credit-risk/26-model-risk-and-validation/validation.csv};
\legend{quarterly steps,monthly steps}
\end{axis}
\end{tikzpicture}

\omcaption{Largest absolute error of the trinomial tree against the closed form, by expiry and step
size, over the validation grid. Refining the step cuts the error by a factor of about three to five, most at
short expiries, where the exercise date falls after only a few steps. Data: the chapter's
tutorial.}\label{fig:rc:model-risk-and-validation:validation}
\end{omfigure}

\section{Inventories, tiering and limitations}

\begin{definition}[Model inventory, model tiering]\label{def:rc:model-risk-and-validation:inventory}
A \emph{model inventory}\index{model inventory} records every model in use with its owner, purpose,
users, validation status and known limitations. \emph{Model tiering}\index{model tiering} classifies
models by the damage they can do (materiality of use, complexity, uncertainty of inputs and
assumptions) and sets the depth and frequency of validation accordingly.
\end{definition}

\begin{example}[The firm's inventory]\label{ex:rc:model-risk-and-validation:inventory}
Scoring eight of this book's models from 1 to 3 on materiality (counted twice), complexity and
uncertainty, and placing tier~1 at a score of 10 or more and tier~2 at 7 or more (an illustrative
policy), gives five tier~1 models (the SABR cube, the \omterm{def:rc:mortgage-modelling:model}{prepayment model}, the exposure engine, the VaR
model and the deposit model), two tier~2 (the curve and the Bermudan pricer) and one tier~3 (the
margin forecast). Tier~1 models are revalidated every year, tier~2 every two, tier~3 every three.
\end{example}

\begin{definition}[Post-model adjustment]\label{def:rc:model-risk-and-validation:pma}
A \emph{post-model adjustment}\index{post-model adjustment} is a documented change to a model's output
made outside the model to correct a known weakness until the model is fixed; it needs an owner, a
rationale, a size and an expiry.
\end{definition}

\section{Benchmarking and outcomes analysis}

Benchmarking finds implementation errors that documentation hides. The 2012 error is a case: the
documentation said relative changes, the spreadsheet computed half of them, and no independent
implementation compared the two.

\begin{example}[Half the volatility]\label{ex:rc:model-risk-and-validation:half}
An illustrative synthetic-credit book sells EUR~1 billion of protection on a 0--3\% equity tranche in the
large-pool model and buys 12.0 times as much index protection as a hedge (hazard 1\%, correlation 20\%).
On a year of illustrative daily moves of the hazard rate and the correlation, its historical 99\% VaR is
EUR~14.06 million with relative changes divided by the average of old and new values, and 7.01 million
when they are divided by the sum: 49.9\% of the intended figure
(\cref{fig:rc:model-risk-and-validation:losses}). A benchmark that recomputed one scenario by hand
would have found it.
\end{example}

\begin{omfigure}
\begin{tikzpicture}
\begin{axis}[width=11.5cm,height=5.2cm,
  xlabel={rank of the scenario loss},ylabel={loss (EUR million)},
  tick label style={font=\small},label style={font=\small},
  xmin=1,xmax=40,ymin=0,ymax=16,grid=major,grid style={gray!25},
  legend columns=2,legend style={font=\footnotesize,at={(0.5,-0.25)},anchor=north,draw=none,fill=none,/tikz/every even column/.append style={column sep=8pt}},
  table/col sep=comma]
\addplot[omDef,very thick,mark=*,mark size=1pt] table[x=rank,y=intended]{figdata/rates-credit-risk/26-model-risk-and-validation/losses.csv};
\addplot[omThm,very thick,mark=square*,mark size=1pt] table[x=rank,y=error]{figdata/rates-credit-risk/26-model-risk-and-validation/losses.csv};
\draw[omExo,dashed] (axis cs:3,0) -- (axis cs:3,16);
\legend{changes divided by the average,changes divided by the sum}
\end{axis}
\end{tikzpicture}

\omcaption{The forty largest scenario losses of the illustrative synthetic-credit book, with relative
changes computed as intended and as in the 2012 spreadsheet. The dashed line marks the scenario that sets
the 99\% VaR: every loss, and the VaR, is about halved. Data: illustrative history; the chapter's
tutorial.}\label{fig:rc:model-risk-and-validation:losses}
\end{omfigure}

\begin{dated}{2026-09}{The 2012 model}\label{dat:rc:model-risk-and-validation:cio}
According to JPMorgan's management task force report (January 2013), the new VaR model for the Chief
Investment Office's synthetic credit portfolio was authorised on 30 January 2012 and ended the breaches
of the firm-wide VaR limit; errors later found in the model included a spreadsheet that divided the
change in a hazard rate or correlation by the sum of the old and new values instead of their average,
likely muting volatility by a factor of two. The report also found that price testing relied on
spreadsheets that had not been vetted.
\end{dated}

\omterm{def:rc:model-risk-and-validation:bench}{Outcomes analysis} catches what benchmarking cannot: a model that is implemented as specified but
wrong about the world. Its tools are the backtests of chapters~21 and~25, the attribution test of
chapter~23 and the monitoring of inputs against the ranges validated.

\section{Tutorial: a validation report}\label{tut:rc:model-risk-and-validation:tutorial}

\begin{tutorial}
\textbf{Goal.} Validate chapter~7's tree against the closed form over a grid, record the failures,
tier the firm's models, and replay the averaging error. \textbf{End state:} the numbers of
\cref{ex:rc:model-risk-and-validation:tree,ex:rc:model-risk-and-validation:inventory,ex:rc:model-risk-and-validation:half}
and the charts, written up in the structure below.
\begin{tutsteps}
\item \textbf{Scope}: the model, its uses, the parameter ranges in use (\texttt{GRID}).
\item \textbf{Benchmark}: \texttt{validation()} runs \texttt{benchmark} with \texttt{tree\_price} against \texttt{closed\_price}.
\item \textbf{Findings}: \texttt{errors\_by\_expiry()}; severity and conditions of use.
\item \textbf{Inventory}: \texttt{inventory()}, \texttt{tiers()}; \texttt{averaging\_error()};
\texttt{fig\_rc\_modelval.py} writes the charts.
\end{tutsteps}
\textbf{What to change next.} Add a Monte Carlo benchmark of the same swaptions; add a finer step
grid near the exercise date and rerun the one-year failures.
\end{tutorial}

\section{Build: model validation tools}\label{bld:rc:model-risk-and-validation:modelval}

\begin{build}
\bfield{Purpose} The firm's inventory and validation tooling, used for every component of this series
before it enters the risk engine.
\bfield{Interface} \texttt{ModelRecord} (with \texttt{tier}, \texttt{revalidate\_years});
\texttt{inventory\_summary}; \texttt{benchmark(candidate, reference, grid, abs\_tol, rel\_tol)},
\texttt{BenchmarkResult}, \texttt{summary}; \texttt{binomial\_tail}; \texttt{relative\_change}.
\bfield{Rules} Every grid point kept, not only failures; tolerances in business units; tiering scores and
intervals are a policy input.
\bfield{Acceptance tests} \texttt{code/firm/modelval/tests/}: tiering boundaries; the harness finds the
single planted failure; binomial tail against the Basel table; the two relative changes differ by a factor
of two.
\bfield{Stretch} Validation reports generated from records; monitoring of input ranges in production;
a findings tracker with owners and deadlines.
\end{build}

\begin{omsources}
\item JPMorgan Chase \& Co., \emph{Report of JPMorgan Chase \& Co. Management Task Force Regarding 2012
CIO Losses}, 16 January 2013.
\item Board of Governors of the Federal Reserve System, OCC and FDIC, \emph{Revised Guidance on Model
Risk Management}, SR 26-2, 17 April 2026.
\item Prudential Regulation Authority, SS1/23, \emph{Model risk management principles for banks}, 2023.
\end{omsources}

\section{Exercises}

\begin{exercise}[$\star$]\label{exo:rc:model-risk-and-validation:1}
A rate moves from 1.00\% to 1.10\%. Give the relative change as intended and as computed in the
2012 spreadsheet.
\end{exercise}

\begin{exercise}[$\star$]\label{exo:rc:model-risk-and-validation:2}
Give the tier of a model scored 3 on materiality, 1 on complexity and 1 on uncertainty under the
chapter's rule.
\end{exercise}

\begin{exercise}[$\star$]\label{exo:rc:model-risk-and-validation:3}
A VaR model has 8 exceptions in 250 days at 99\%. How likely is that for a correct model?
\end{exercise}

\begin{exercise}[$\star\star$]\label{exo:rc:model-risk-and-validation:4}
Why would a spreadsheet error that halves volatility not be caught by a VaR backtest within weeks?
\end{exercise}

\begin{exercise}[$\star\star$]\label{exo:rc:model-risk-and-validation:5}
Why is a closed form of the same model a weaker benchmark than an alternative model, and why is it
still useful?
\end{exercise}

\begin{exercise}[$\star\star$]\label{exo:rc:model-risk-and-validation:6}
Under SR 26-2's definition, is the 2012 spreadsheet a model? What follows for its control?
\end{exercise}

\begin{exercise}[$\star\star\star$]\label{exo:rc:model-risk-and-validation:7}
\emph{Coding.} Rerun the tree validation with steps of one twenty-fourth of a year for one-year
expiries only, and report the failures.
\end{exercise}

\begin{exercise}[$\star\star\star$]\label{exo:rc:model-risk-and-validation:8}
\emph{Find the flaw.} ``The new model was approved by the review group, so the lower VaR is the right
number.''
\end{exercise}

\section{Problem: The Averaging Error}

\begin{problem}[{Weekend problem --- a factor of two}]\label{pb:rc:model-risk-and-validation:1}
The illustrative synthetic-credit book of \cref{ex:rc:model-risk-and-validation:half} is measured by a
VaR model implemented in a spreadsheet that divides relative changes by the sum of old and new values.

\medskip\noindent\textbf{Part I --- The error.}
\begin{enumerate}
\item Show that the error halves every relative change.
\item Give the VaR with and without the error, and the ratio.
\item Why is the ratio not exactly one half?
\item Give the hedge ratio of the book and why it is so large.
\item What limit behaviour would the error produce?
\end{enumerate}

\medskip\noindent\textbf{Part II --- Detection.}
\begin{enumerate}[resume]
\item Which validation step would have caught it, and how quickly?
\item What would \omterm{def:rc:model-risk-and-validation:bench}{outcomes analysis} have shown, and when?
\item What role did the timing of the model change play in 2012?
\item Why did the model's documentation not reveal the error?
\item What control on spreadsheets would you require?
\end{enumerate}

\medskip\noindent\textbf{Part III --- Governance.}
\begin{enumerate}[resume]
\item Who should approve a model that lowers a desk's VaR during a limit breach?
\item What conditions of use would you attach?
\item How should limits be treated when a model changes?
\item Which tier would you give this model, and why?
\item What does \omterm{def:rc:model-risk-and-validation:challenge}{effective challenge} require of the validator here?
\end{enumerate}

\medskip\noindent\textbf{Part IV --- Judgement.}
\begin{enumerate}[resume]
\item Was the 2012 loss caused by the model?
\item What does SR 26-2's exclusion of simple spreadsheet arithmetic mean for cases like this?
\item How would you design a validation that catches implementation errors cheaply?
\item State the \emph{named result}: the VaR reduction produced by dividing a rate change by the sum of
two rates instead of their average.
\item In one sentence: what is \omterm{def:rc:model-risk-and-validation:risk}{model risk}?
\end{enumerate}
\end{problem}

\section{Interview questions}

\begin{interviewq}[$\star$ \iqroles{risk, bank}]\label{iq:rc:model-risk-and-validation:1}
What does a \omterm{def:rc:model-risk-and-validation:validation}{model validation} report contain?
\end{interviewq}

\begin{interviewq}[$\star\star$ \iqroles{researcher, risk}]\label{iq:rc:model-risk-and-validation:2}
How would you validate a pricing model for \omterm{def:rc:bermudans-and-callables:bermudan}{Bermudan swaptions}?
\end{interviewq}

\begin{interviewq}[$\star\star$ \iqroles{developer}]\label{iq:rc:model-risk-and-validation:3}
How do you test a numerical pricer so that implementation errors are found early?
\end{interviewq}

\begin{interviewq}[$\star\star$ \iqroles{risk}]\label{iq:rc:model-risk-and-validation:4}
What is \omterm{def:rc:model-risk-and-validation:challenge}{effective challenge}, and how do you know it is happening?
\end{interviewq}

\begin{interviewq}[$\star\star\star$ \iqroles{bank, risk}]\label{iq:rc:model-risk-and-validation:5}
What went wrong with the 2012 CIO VaR model, and which controls would have stopped it?
\end{interviewq}

\begin{interviewq}[$\star\star\star$ \iqroles{researcher}]\label{iq:rc:model-risk-and-validation:6}
How would you quantify \omterm{def:rc:model-risk-and-validation:risk}{model risk} in a price, for a reserve?
\end{interviewq}
